
\documentclass{article}
\usepackage{colortbl}
\usepackage{makecell}
\usepackage{multirow}
\usepackage{supertabular}

\begin{document}

\newcounter{utterance}

\centering \large Interaction Transcript for game `clean\_up', experiment `1\_medium\_3obj\_uk', episode 0 with gemma{-}4{-}26b{-}a4b{-}it.
\vspace{24pt}

{ \footnotesize  \setcounter{utterance}{1}
\setlength{\tabcolsep}{0pt}
\begin{supertabular}{c@{$\;$}|p{.15\linewidth}@{}p{.15\linewidth}p{.15\linewidth}p{.15\linewidth}p{.15\linewidth}p{.15\linewidth}}
    \# & $\;$A & \multicolumn{4}{c}{Game Master} & $\;\:$B\\
    \hline

    \theutterance \stepcounter{utterance}  
    & & \multicolumn{4}{p{0.6\linewidth}}{
        \cellcolor[rgb]{0.9,0.9,0.9}{
            \makecell[{{p{\linewidth}}}]{
                \texttt{\tiny{[P1$\langle$GM]}}
                \texttt{Я ваш ведучий гри, і ви граєте у спільну гру з наступною сіткою як ігровим полем:} \\
\\ 
\texttt{```} \\
\texttt{ 1234567} \\
\texttt{╔════╤══╗} \\
\texttt{║◌C◌◌│◌◌║ 1} \\
\texttt{║◌◌L◌├──╢ 2} \\
\texttt{║◌◌◌P│◌◌║ 3} \\
\texttt{║◌┌──┼──╢ 4} \\
\texttt{║◌│◌◌│◌◌║ 5} \\
\texttt{╟─┴┬─┤◌◌║ 6} \\
\texttt{║◌◌│◌│◌◌║ 7} \\
\texttt{╚══╧═╧══╝} \\
\\ 
\texttt{```} \\
\\ 
\texttt{* Верхній край показує координати x, що зростають праворуч, а правий край — координати y, що зростають донизу.} \\
\texttt{* Наступні об'єкти випадково розміщені на вашій сітці: 'C', 'L', 'P'.} \\
\\ 
\\ 
\texttt{Інший гравець бачить варіант ігрового поля, де об'єкти розміщені в інших випадкових місцях. Ви не бачите поле іншого гравця, а він не бачить вашого.} \\
\\ 
\\ 
\texttt{**Мета:**} \\
\\ 
\texttt{Обидва гравці повинні переміщати об'єкти на своєму фоні так, щоб однакові об'єкти опинилися в однакових координатах. Ви повинні спілкуватися з іншим гравцем, щоб узгодити спільну цільову конфігурацію.} \\
\\ 
\\ 
\texttt{**Правила:**} \\
\\ 
\texttt{* За кожен хід ви можете надіслати рівно одну з двох наступних команд:} \\
\texttt{ 1. `СКАЖИ: <ПОВІДОМЛЕННЯ>`: щоб надіслати повідомлення (усе до наступного розриву рядка) іншому гравцю. Я передам його вашому партнеру.} \\
\texttt{2. `ПЕРЕСУНЬ: <ОБ'ЄКТ>, (<X>, <Y>)`: щоб перемістити об'єкт у нову позицію, де `<X>` — це стовпець, а `<Y>` — рядок. Я повідомлю вас, чи був ваш хід успішним.} \\
\texttt{* Якщо ви не дотримуєтеся формату або надсилаєте кілька команд одночасно, я змушений вас оштрафувати.} \\
\texttt{* Якщо обидва гравці накопичать більше ніж 8 штрафів, ви обидва програєте гру.} \\
\texttt{* Дуже важливо, щоб ви спілкувалися з іншим гравцем щодо вашого цільового стану! *Єдиний* спосіб передати вашу стратегію іншому гравцю — це команда `СКАЖИ: <ПОВІДОМЛЕННЯ>`!} \\
\\ 
\\ 
\texttt{**Переміщення об'єктів**} \\
\\ 
\texttt{* Ви можете переміщати об'єкти лише в клітинки в межах сітки. Цільова клітинка повинна бути порожньою, тобто містити лише символ '◌'.} \\
\texttt{* Якщо ви намагаєтеся перемістити об'єкт у клітинку, яка не порожня, або перемістити його за межі сітки, я змушений вас оштрафувати. Ви отримаєте ще одну спробу.} \\
\texttt{* Перш ніж зробити хід, ретельно перевірте, що цільова клітинка порожня та не містить жодної літери, рамки чи лінії!} \\
\\ 
\\ 
\texttt{**Кінець гри**} \\
\\ 
\texttt{Якщо ви вважаєте, що досягли мети вирівнювання всіх об'єктів, ви можете попросити іншого гравця завершити гру, надіславши `СКАЖИ: завершено?`. Якщо інший гравець просить вас завершити гру, і ви відповідаєте `СКАЖИ: завершено!`, гра завершиться.} \\
\\ 
\texttt{Обидва гравці перемагають, якщо гра завершується протягом 12 раундів, де один раунд визначається як надсилання обома гравцями по одній дійсній команді.} \\
\\ 
\texttt{**Підрахунок очок:**} \\
\\ 
\texttt{Чим ближче однакові об'єкти на обох ігрових полях, тим більше очок ви отримуєте. Штрафи зменшують ваші очки. Чи зможете ви побити рекорд?} \\
\\ 
\\ 
\texttt{Будь ласка, надішліть повідомлення іншому гравцю, щоб почати гру!} \\
            }
        }
    }
    & \\ \\

    \theutterance \stepcounter{utterance}  
    & \multicolumn{4}{p{0.6\linewidth}}{
        \cellcolor[rgb]{0.9,0.9,0.9}{
            \makecell[{{p{\linewidth}}}]{
                \texttt{\tiny{[P1$\rangle$GM]}}
                \texttt{СКАЖИ: Привіт! Я бачу свої об'єкти на позиціях: C(2,1), L(3,2), P(4,3). Який стан твого поля? Давай узгодимо спільні координати для C, L та P.} \\
            }
        }
    }
    & & \\ \\

    \theutterance \stepcounter{utterance}  
    & & & \multicolumn{2}{p{0.3\linewidth}}{
        \cellcolor[rgb]{0.9,0.9,0.9}{
            \makecell[{{p{\linewidth}}}]{
                \texttt{\tiny{[GM$|$GM]}}
                \texttt{СКАЖИ: Привіт! Я бачу свої об'єкти на позиціях: C(2,1), L(3,2), P(4,3). Який стан твого поля? Давай узгодимо спільні координати для C, L та P.} \\
            }
        }
    }
    & & \\ \\

    \theutterance \stepcounter{utterance}  
    & & \multicolumn{4}{p{0.6\linewidth}}{
        \cellcolor[rgb]{0.9,0.9,0.9}{
            \makecell[{{p{\linewidth}}}]{
                \texttt{\tiny{[GM$\rangle$P2]}}
                \texttt{Я ваш ведучий гри, і ви граєте у спільну гру з наступною сіткою як ігровим полем:} \\
\\ 
\texttt{```} \\
\texttt{ 1234567} \\
\texttt{╔════╤══╗} \\
\texttt{║◌◌P◌│◌◌║ 1} \\
\texttt{║◌◌C◌├──╢ 2} \\
\texttt{║◌◌◌◌│◌◌║ 3} \\
\texttt{║◌┌──┼──╢ 4} \\
\texttt{║◌│◌◌│◌◌║ 5} \\
\texttt{╟─┴┬─┤◌◌║ 6} \\
\texttt{║L◌│◌│◌◌║ 7} \\
\texttt{╚══╧═╧══╝} \\
\\ 
\texttt{```} \\
\\ 
\texttt{* Верхній край показує координати x, що зростають праворуч, а правий край — координати y, що зростають донизу.} \\
\texttt{* Наступні об'єкти випадково розміщені на вашій сітці: 'C', 'L', 'P'.} \\
\\ 
\\ 
\texttt{Інший гравець бачить варіант ігрового поля, де об'єкти розміщені в інших випадкових місцях. Ви не бачите поле іншого гравця, а він не бачить вашого.} \\
\\ 
\\ 
\texttt{**Мета:**} \\
\\ 
\texttt{Обидва гравці повинні переміщати об'єкти на своєму фоні так, щоб однакові об'єкти опинилися в однакових координатах. Ви повинні спілкуватися з іншим гравцем, щоб узгодити спільну цільову конфігурацію.} \\
\\ 
\\ 
\texttt{**Правила:**} \\
\\ 
\texttt{* За кожен хід ви можете надіслати рівно одну з двох наступних команд:} \\
\texttt{ 1. `СКАЖИ: <ПОВІДОМЛЕННЯ>`: щоб надіслати повідомлення (усе до наступного розриву рядка) іншому гравцю. Я передам його вашому партнеру.} \\
\texttt{2. `ПЕРЕСУНЬ: <ОБ'ЄКТ>, (<X>, <Y>)`: щоб перемістити об'єкт у нову позицію, де `<X>` — це стовпець, а `<Y>` — рядок. Я повідомлю вас, чи був ваш хід успішним.} \\
\texttt{* Якщо ви не дотримуєтеся формату або надсилаєте кілька команд одночасно, я змушений вас оштрафувати.} \\
\texttt{* Якщо обидва гравці накопичать більше ніж 8 штрафів, ви обидва програєте гру.} \\
\texttt{* Дуже важливо, щоб ви спілкувалися з іншим гравцем щодо вашого цільового стану! *Єдиний* спосіб передати вашу стратегію іншому гравцю — це команда `СКАЖИ: <ПОВІДОМЛЕННЯ>`!} \\
\\ 
\\ 
\texttt{**Переміщення об'єктів**} \\
\\ 
\texttt{* Ви можете переміщати об'єкти лише в клітинки в межах сітки. Цільова клітинка повинна бути порожньою, тобто містити лише символ '◌'.} \\
\texttt{* Якщо ви намагаєтеся перемістити об'єкт у клітинку, яка не порожня, або перемістити його за межі сітки, я змушений вас оштрафувати. Ви отримаєте ще одну спробу.} \\
\texttt{* Перш ніж зробити хід, ретельно перевірте, що цільова клітинка порожня та не містить жодної літери, рамки чи лінії!} \\
\\ 
\\ 
\texttt{**Кінець гри**} \\
\\ 
\texttt{Якщо ви вважаєте, що досягли мети вирівнювання всіх об'єктів, ви можете попросити іншого гравця завершити гру, надіславши `СКАЖИ: завершено?`. Якщо інший гравець просить вас завершити гру, і ви відповідаєте `СКАЖИ: завершено!`, гра завершиться.} \\
\\ 
\texttt{Обидва гравці перемагають, якщо гра завершується протягом 12 раундів, де один раунд визначається як надсилання обома гравцями по одній дійсній команді.} \\
\\ 
\texttt{**Підрахунок очок:**} \\
\\ 
\texttt{Чим ближче однакові об'єкти на обох ігрових полях, тим більше очок ви отримуєте. Штрафи зменшують ваші очки. Чи зможете ви побити рекорд?} \\
\\ 
\texttt{Інший гравець почав гру, надіславши це повідомлення:} \\
\\ 
\texttt{"Привіт! Я бачу свої об'єкти на позиціях: C(2,1), L(3,2), P(4,3). Який стан твого поля? Давай узгодимо спільні координати для C, L та P."} \\
\\ 
\texttt{Яка ваша перша команда?} \\
            }
        }
    }
    & \\ \\

    \theutterance \stepcounter{utterance}  
    & & & \multicolumn{4}{p{0.6\linewidth}}{
        \cellcolor[rgb]{0.9,0.9,0.9}{
            \makecell[{{p{\linewidth}}}]{
                \texttt{\tiny{[GM$\langle$P2]}}
                \texttt{СКАЖИ: Привіт! Мої об'єкти на позиціях: P(3,1), C(3,2), L(1,7). Давай спробуємо звести їх до твоїх координат: C(2,1), L(3,2), P(4,3). Чи згодна ти на таку ціль?} \\
            }
        }
    }
     \\ \\

    \theutterance \stepcounter{utterance}  
    & & & \multicolumn{2}{p{0.3\linewidth}}{
        \cellcolor[rgb]{0.9,0.9,0.9}{
            \makecell[{{p{\linewidth}}}]{
                \texttt{\tiny{[GM$|$GM]}}
                \texttt{СКАЖИ: Привіт! Мої об'єкти на позиціях: P(3,1), C(3,2), L(1,7). Давай спробуємо звести їх до твоїх координат: C(2,1), L(3,2), P(4,3). Чи згодна ти на таку ціль?} \\
            }
        }
    }
    & & \\ \\

    \theutterance \stepcounter{utterance}  
    & & \multicolumn{4}{p{0.6\linewidth}}{
        \cellcolor[rgb]{0.9,0.9,0.9}{
            \makecell[{{p{\linewidth}}}]{
                \texttt{\tiny{[P1$\langle$GM]}}
                \texttt{Ваше повідомлення передано іншому гравцю.} \\
\\ 
\\ 
\texttt{Зараз ви граєте раунд 1 з максимум 12.} \\
\texttt{Разом ви накопичили 0 з 8 штрафів.} \\
\texttt{Інший гравець надіслав це повідомлення:} \\
\\ 
\\ 
\texttt{"Привіт! Мої об'єкти на позиціях: P(3,1), C(3,2), L(1,7). Давай спробуємо звести їх до твоїх координат: C(2,1), L(3,2), P(4,3). Чи згодна ти на таку ціль?"} \\
\\ 
\\ 
\texttt{Яка ваша наступна команда?} \\
            }
        }
    }
    & \\ \\

    \theutterance \stepcounter{utterance}  
    & \multicolumn{4}{p{0.6\linewidth}}{
        \cellcolor[rgb]{0.9,0.9,0.9}{
            \makecell[{{p{\linewidth}}}]{
                \texttt{\tiny{[P1$\rangle$GM]}}
                \texttt{СКАЖИ: Так, я згодна! Давай тримати ціль: C(2,1), L(3,2), P(4,3). Починаймо переміщення.} \\
            }
        }
    }
    & & \\ \\

    \theutterance \stepcounter{utterance}  
    & & & \multicolumn{2}{p{0.3\linewidth}}{
        \cellcolor[rgb]{0.9,0.9,0.9}{
            \makecell[{{p{\linewidth}}}]{
                \texttt{\tiny{[GM$|$GM]}}
                \texttt{СКАЖИ: Так, я згодна! Давай тримати ціль: C(2,1), L(3,2), P(4,3). Починаймо переміщення.} \\
            }
        }
    }
    & & \\ \\

    \theutterance \stepcounter{utterance}  
    & & \multicolumn{4}{p{0.6\linewidth}}{
        \cellcolor[rgb]{0.9,0.9,0.9}{
            \makecell[{{p{\linewidth}}}]{
                \texttt{\tiny{[GM$\rangle$P2]}}
                \texttt{Ваше повідомлення передано іншому гравцю.} \\
\\ 
\\ 
\texttt{Зараз ви граєте раунд 2 з максимум 12.} \\
\texttt{Разом ви накопичили 0 з 8 штрафів.} \\
\texttt{Інший гравець надіслав це повідомлення:} \\
\\ 
\\ 
\texttt{"Так, я згодна! Давай тримати ціль: C(2,1), L(3,2), P(4,3). Починаймо переміщення."} \\
\\ 
\\ 
\texttt{Яка ваша наступна команда?} \\
            }
        }
    }
    & \\ \\

    \theutterance \stepcounter{utterance}  
    & & & \multicolumn{4}{p{0.6\linewidth}}{
        \cellcolor[rgb]{0.9,0.9,0.9}{
            \makecell[{{p{\linewidth}}}]{
                \texttt{\tiny{[GM$\langle$P2]}}
                \texttt{ПЕРЕСУНЬ: C, (2, 2)} \\
            }
        }
    }
     \\ \\

    \theutterance \stepcounter{utterance}  
    & & & \multicolumn{2}{p{0.3\linewidth}}{
        \cellcolor[rgb]{0.9,0.9,0.9}{
            \makecell[{{p{\linewidth}}}]{
                \texttt{\tiny{[GM$|$GM]}}
                \texttt{ПЕРЕСУНЬ: C, (2, 2)} \\
            }
        }
    }
    & & \\ \\

    \theutterance \stepcounter{utterance}  
    & & & \multicolumn{2}{p{0.3\linewidth}}{
        \cellcolor[rgb]{0.9,0.9,0.9}{
            \makecell[{{p{\linewidth}}}]{
                \texttt{\tiny{[GM$|$GM]}}
                \texttt{C moved to (2, 2)} \\
            }
        }
    }
    & & \\ \\

    \theutterance \stepcounter{utterance}  
    & & \multicolumn{4}{p{0.6\linewidth}}{
        \cellcolor[rgb]{0.9,0.9,0.9}{
            \makecell[{{p{\linewidth}}}]{
                \texttt{\tiny{[P1$\langle$GM]}}
                \texttt{Ваше повідомлення передано іншому гравцю.} \\
\\ 
\\ 
\texttt{Зараз ви граєте раунд 2 з максимум 12.} \\
\texttt{Разом ви накопичили 0 з 8 штрафів.} \\
\texttt{Інший гравець перемістив об'єкт на своїй сітці.} \\
\\ 
\\ 
\texttt{Яка ваша наступна команда?} \\
            }
        }
    }
    & \\ \\

    \theutterance \stepcounter{utterance}  
    & \multicolumn{4}{p{0.6\linewidth}}{
        \cellcolor[rgb]{0.9,0.9,0.9}{
            \makecell[{{p{\linewidth}}}]{
                \texttt{\tiny{[P1$\rangle$GM]}}
                \texttt{СКАЖИ: Який твій наступний хід? Я збираюся перемістити P(4,3) на (4,3), але воно вже там. Спробуймо спочатку змістити щось інше, якщо потрібно. Який твій статус?} \\
            }
        }
    }
    & & \\ \\

    \theutterance \stepcounter{utterance}  
    & & & \multicolumn{2}{p{0.3\linewidth}}{
        \cellcolor[rgb]{0.9,0.9,0.9}{
            \makecell[{{p{\linewidth}}}]{
                \texttt{\tiny{[GM$|$GM]}}
                \texttt{СКАЖИ: Який твій наступний хід? Я збираюся перемістити P(4,3) на (4,3), але воно вже там. Спробуймо спочатку змістити щось інше, якщо потрібно. Який твій статус?} \\
            }
        }
    }
    & & \\ \\

    \theutterance \stepcounter{utterance}  
    & & \multicolumn{4}{p{0.6\linewidth}}{
        \cellcolor[rgb]{0.9,0.9,0.9}{
            \makecell[{{p{\linewidth}}}]{
                \texttt{\tiny{[GM$\rangle$P2]}}
                \texttt{Об'єкт 'C' успішно переміщено в (2,2). Ваша оновлена сітка виглядає так:} \\
\texttt{```} \\
\texttt{ 1234567} \\
\texttt{╔════╤══╗} \\
\texttt{║◌◌P◌│◌◌║ 1} \\
\texttt{║◌C◌◌├──╢ 2} \\
\texttt{║◌◌◌◌│◌◌║ 3} \\
\texttt{║◌┌──┼──╢ 4} \\
\texttt{║◌│◌◌│◌◌║ 5} \\
\texttt{╟─┴┬─┤◌◌║ 6} \\
\texttt{║L◌│◌│◌◌║ 7} \\
\texttt{╚══╧═╧══╝} \\
\texttt{```} \\
\texttt{Зараз ви граєте раунд 3 з максимум 12.} \\
\texttt{Разом ви накопичили 0 з 8 штрафів.} \\
\texttt{Інший гравець надіслав це повідомлення:} \\
\\ 
\\ 
\texttt{"Який твій наступний хід? Я збираюся перемістити P(4,3) на (4,3), але воно вже там. Спробуймо спочатку змістити щось інше, якщо потрібно. Який твій статус?"} \\
\\ 
\\ 
\texttt{Яка ваша наступна команда?} \\
            }
        }
    }
    & \\ \\

    \theutterance \stepcounter{utterance}  
    & & & \multicolumn{4}{p{0.6\linewidth}}{
        \cellcolor[rgb]{0.9,0.9,0.9}{
            \makecell[{{p{\linewidth}}}]{
                \texttt{\tiny{[GM$\langle$P2]}}
                \texttt{ПЕРЕСУНЬ: P, (4, 1)} \\
            }
        }
    }
     \\ \\

    \theutterance \stepcounter{utterance}  
    & & & \multicolumn{2}{p{0.3\linewidth}}{
        \cellcolor[rgb]{0.9,0.9,0.9}{
            \makecell[{{p{\linewidth}}}]{
                \texttt{\tiny{[GM$|$GM]}}
                \texttt{ПЕРЕСУНЬ: P, (4, 1)} \\
            }
        }
    }
    & & \\ \\

    \theutterance \stepcounter{utterance}  
    & & & \multicolumn{2}{p{0.3\linewidth}}{
        \cellcolor[rgb]{0.9,0.9,0.9}{
            \makecell[{{p{\linewidth}}}]{
                \texttt{\tiny{[GM$|$GM]}}
                \texttt{P moved to (4, 1)} \\
            }
        }
    }
    & & \\ \\

    \theutterance \stepcounter{utterance}  
    & & \multicolumn{4}{p{0.6\linewidth}}{
        \cellcolor[rgb]{0.9,0.9,0.9}{
            \makecell[{{p{\linewidth}}}]{
                \texttt{\tiny{[P1$\langle$GM]}}
                \texttt{Ваше повідомлення передано іншому гравцю.} \\
\\ 
\\ 
\texttt{Зараз ви граєте раунд 3 з максимум 12.} \\
\texttt{Разом ви накопичили 0 з 8 штрафів.} \\
\texttt{Інший гравець перемістив об'єкт на своїй сітці.} \\
\\ 
\\ 
\texttt{Яка ваша наступна команда?} \\
            }
        }
    }
    & \\ \\

    \theutterance \stepcounter{utterance}  
    & \multicolumn{4}{p{0.6\linewidth}}{
        \cellcolor[rgb]{0.9,0.9,0.9}{
            \makecell[{{p{\linewidth}}}]{
                \texttt{\tiny{[P1$\rangle$GM]}}
                \texttt{СКАЖИ: Я бачу, що ти робиш ходи, але я не знаю, які саме. Будь ласка, напиши, який об'єкт і куди ти перемістив, щоб ми могли синхронізуватися. Моя ціль залишається: C(2,1), L(3,2), P(4,3).} \\
            }
        }
    }
    & & \\ \\

    \theutterance \stepcounter{utterance}  
    & & & \multicolumn{2}{p{0.3\linewidth}}{
        \cellcolor[rgb]{0.9,0.9,0.9}{
            \makecell[{{p{\linewidth}}}]{
                \texttt{\tiny{[GM$|$GM]}}
                \texttt{СКАЖИ: Я бачу, що ти робиш ходи, але я не знаю, які саме. Будь ласка, напиши, який об'єкт і куди ти перемістив, щоб ми могли синхронізуватися. Моя ціль залишається: C(2,1), L(3,2), P(4,3).} \\
            }
        }
    }
    & & \\ \\

    \theutterance \stepcounter{utterance}  
    & & \multicolumn{4}{p{0.6\linewidth}}{
        \cellcolor[rgb]{0.9,0.9,0.9}{
            \makecell[{{p{\linewidth}}}]{
                \texttt{\tiny{[GM$\rangle$P2]}}
                \texttt{Об'єкт 'P' успішно переміщено в (4,1). Ваша оновлена сітка виглядає так:} \\
\texttt{```} \\
\texttt{ 1234567} \\
\texttt{╔════╤══╗} \\
\texttt{║◌◌◌P│◌◌║ 1} \\
\texttt{║◌C◌◌├──╢ 2} \\
\texttt{║◌◌◌◌│◌◌║ 3} \\
\texttt{║◌┌──┼──╢ 4} \\
\texttt{║◌│◌◌│◌◌║ 5} \\
\texttt{╟─┴┬─┤◌◌║ 6} \\
\texttt{║L◌│◌│◌◌║ 7} \\
\texttt{╚══╧═╧══╝} \\
\texttt{```} \\
\texttt{Зараз ви граєте раунд 4 з максимум 12.} \\
\texttt{Разом ви накопичили 0 з 8 штрафів.} \\
\texttt{Інший гравець надіслав це повідомлення:} \\
\\ 
\\ 
\texttt{"Я бачу, що ти робиш ходи, але я не знаю, які саме. Будь ласка, напиши, який об'єкт і куди ти перемістив, щоб ми могли синхронізуватися. Моя ціль залишається: C(2,1), L(3,2), P(4,3)."} \\
\\ 
\\ 
\texttt{Яка ваша наступна команда?} \\
            }
        }
    }
    & \\ \\

    \theutterance \stepcounter{utterance}  
    & & & \multicolumn{4}{p{0.6\linewidth}}{
        \cellcolor[rgb]{0.9,0.9,0.9}{
            \makecell[{{p{\linewidth}}}]{
                \texttt{\tiny{[GM$\langle$P2]}}
                \texttt{СКАЖИ: Зрозумів! Я перемістив C у (2,2) та P у (4,1). Моя ціль: C(2,1), L(3,2), P(4,3). Зараз я буду переміщувати об'єкти до цих координат. Наступним планую перемістити P у (4,3).} \\
            }
        }
    }
     \\ \\

    \theutterance \stepcounter{utterance}  
    & & & \multicolumn{2}{p{0.3\linewidth}}{
        \cellcolor[rgb]{0.9,0.9,0.9}{
            \makecell[{{p{\linewidth}}}]{
                \texttt{\tiny{[GM$|$GM]}}
                \texttt{СКАЖИ: Зрозумів! Я перемістив C у (2,2) та P у (4,1). Моя ціль: C(2,1), L(3,2), P(4,3). Зараз я буду переміщувати об'єкти до цих координат. Наступним планую перемістити P у (4,3).} \\
            }
        }
    }
    & & \\ \\

    \theutterance \stepcounter{utterance}  
    & & \multicolumn{4}{p{0.6\linewidth}}{
        \cellcolor[rgb]{0.9,0.9,0.9}{
            \makecell[{{p{\linewidth}}}]{
                \texttt{\tiny{[P1$\langle$GM]}}
                \texttt{Ваше повідомлення передано іншому гравцю.} \\
\\ 
\\ 
\texttt{Зараз ви граєте раунд 4 з максимум 12.} \\
\texttt{Разом ви накопичили 0 з 8 штрафів.} \\
\texttt{Інший гравець надіслав це повідомлення:} \\
\\ 
\\ 
\texttt{"Зрозумів! Я перемістив C у (2,2) та P у (4,1). Моя ціль: C(2,1), L(3,2), P(4,3). Зараз я буду переміщувати об'єкти до цих координат. Наступним планую перемістити P у (4,3)."} \\
\\ 
\\ 
\texttt{Яка ваша наступна команда?} \\
            }
        }
    }
    & \\ \\

    \theutterance \stepcounter{utterance}  
    & \multicolumn{4}{p{0.6\linewidth}}{
        \cellcolor[rgb]{0.9,0.9,0.9}{
            \makecell[{{p{\linewidth}}}]{
                \texttt{\tiny{[P1$\rangle$GM]}}
                \texttt{ПЕРЕСУНЬ: C, (2,1)} \\
            }
        }
    }
    & & \\ \\

    \theutterance \stepcounter{utterance}  
    & & & \multicolumn{2}{p{0.3\linewidth}}{
        \cellcolor[rgb]{0.9,0.9,0.9}{
            \makecell[{{p{\linewidth}}}]{
                \texttt{\tiny{[GM$|$GM]}}
                \texttt{ПЕРЕСУНЬ: C, (2,1)} \\
            }
        }
    }
    & & \\ \\

    \theutterance \stepcounter{utterance}  
    & & & \multicolumn{2}{p{0.3\linewidth}}{
        \cellcolor[rgb]{0.9,0.9,0.9}{
            \makecell[{{p{\linewidth}}}]{
                \texttt{\tiny{[GM$|$GM]}}
                \texttt{Штраф: (2,1) не порожня, а містить 'C'.} \\
\\ 
\texttt{Разом ви накопичили 1 з 8 штрафів.} \\
\texttt{Будь ласка, спробуйте ще раз!} \\
            }
        }
    }
    & & \\ \\

    \theutterance \stepcounter{utterance}  
    & & \multicolumn{4}{p{0.6\linewidth}}{
        \cellcolor[rgb]{0.9,0.9,0.9}{
            \makecell[{{p{\linewidth}}}]{
                \texttt{\tiny{[P1$\langle$GM]}}
                \texttt{Штраф: (2,1) не порожня, а містить 'C'.} \\
\\ 
\texttt{Разом ви накопичили 1 з 8 штрафів.} \\
\texttt{Будь ласка, спробуйте ще раз!} \\
            }
        }
    }
    & \\ \\

    \theutterance \stepcounter{utterance}  
    & \multicolumn{4}{p{0.6\linewidth}}{
        \cellcolor[rgb]{0.9,0.9,0.9}{
            \makecell[{{p{\linewidth}}}]{
                \texttt{\tiny{[P1$\rangle$GM]}}
                \texttt{ПЕРЕСУНЬ: L, (3,2)} \\
            }
        }
    }
    & & \\ \\

    \theutterance \stepcounter{utterance}  
    & & & \multicolumn{2}{p{0.3\linewidth}}{
        \cellcolor[rgb]{0.9,0.9,0.9}{
            \makecell[{{p{\linewidth}}}]{
                \texttt{\tiny{[GM$|$GM]}}
                \texttt{ПЕРЕСУНЬ: L, (3,2)} \\
            }
        }
    }
    & & \\ \\

    \theutterance \stepcounter{utterance}  
    & & & \multicolumn{2}{p{0.3\linewidth}}{
        \cellcolor[rgb]{0.9,0.9,0.9}{
            \makecell[{{p{\linewidth}}}]{
                \texttt{\tiny{[GM$|$GM]}}
                \texttt{Штраф: (3,2) не порожня, а містить 'L'.} \\
\\ 
\texttt{Разом ви накопичили 2 з 8 штрафів.} \\
\texttt{Будь ласка, спробуйте ще раз!} \\
            }
        }
    }
    & & \\ \\

    \theutterance \stepcounter{utterance}  
    & & \multicolumn{4}{p{0.6\linewidth}}{
        \cellcolor[rgb]{0.9,0.9,0.9}{
            \makecell[{{p{\linewidth}}}]{
                \texttt{\tiny{[P1$\langle$GM]}}
                \texttt{Штраф: (3,2) не порожня, а містить 'L'.} \\
\\ 
\texttt{Разом ви накопичили 2 з 8 штрафів.} \\
\texttt{Будь ласка, спробуйте ще раз!} \\
            }
        }
    }
    & \\ \\

    \theutterance \stepcounter{utterance}  
    & \multicolumn{4}{p{0.6\linewidth}}{
        \cellcolor[rgb]{0.9,0.9,0.9}{
            \makecell[{{p{\linewidth}}}]{
                \texttt{\tiny{[P1$\rangle$GM]}}
                \texttt{Ой, я припустився помилки. Я намагався перемістити об'єкти в ті самі клітинки, де вони вже знаходяться. Оскільки мої об'єкти вже на цільових позиціях (C на 2,1; L на 3,2; P на 4,3), мені не потрібно їх рухати.} \\
\\ 
\texttt{СКАЖИ: Вибач, я помилився з командами. Мої об'єкти вже знаходяться на цільових позиціях: C(2,1), L(3,2), P(4,3). Якщо твої об'єкти також будуть на цих місцях, ми зможемо завершити гру. Чекаю на твій хід!} \\
            }
        }
    }
    & & \\ \\

    \theutterance \stepcounter{utterance}  
    & & & \multicolumn{2}{p{0.3\linewidth}}{
        \cellcolor[rgb]{0.9,0.9,0.9}{
            \makecell[{{p{\linewidth}}}]{
                \texttt{\tiny{[GM$|$GM]}}
                \texttt{Ой, я припустився помилки. Я намагався перемістити об'єкти в ті самі клітинки, де вони вже знаходяться. Оскільки мої об'єкти вже на цільових позиціях (C на 2,1; L на 3,2; P на 4,3), мені не потрібно їх рухати.} \\
\\ 
\texttt{СКАЖИ: Вибач, я помилився з командами. Мої об'єкти вже знаходяться на цільових позиціях: C(2,1), L(3,2), P(4,3). Якщо твої об'єкти також будуть на цих місцях, ми зможемо завершити гру. Чекаю на твій хід!} \\
            }
        }
    }
    & & \\ \\

    \theutterance \stepcounter{utterance}  
    & & & \multicolumn{2}{p{0.3\linewidth}}{
        \cellcolor[rgb]{0.9,0.9,0.9}{
            \makecell[{{p{\linewidth}}}]{
                \texttt{\tiny{[GM$|$GM]}}
                \texttt{Invalid format: head is not empty: 'Ой, я припустився помилки. Я намагався перемістити об'єкти в ті самі клітинки, де вони вже знаходяться. Оскільки мої об'єкти вже на цільових позиціях (C на 2,1; L на 3,2; P на 4,3), мені не потрібно їх рухати.} \\
\\ 
\texttt{'} \\
            }
        }
    }
    & & \\ \\

    \theutterance \stepcounter{utterance}  
    & & \multicolumn{4}{p{0.6\linewidth}}{
        \cellcolor[rgb]{0.9,0.9,0.9}{
            \makecell[{{p{\linewidth}}}]{
                \texttt{\tiny{[P1$\langle$GM]}}
                \texttt{Штраф: Ваше повідомлення не повинно містити нічого перед командою!} \\
\\ 
\texttt{Переконайтеся, що ваша відповідь містить лише `СКАЖИ: <ПОВІДОМЛЕННЯ>` або `ПЕРЕСУНЬ: <ОБ'ЄКТ>, (<X>, <Y>)` і нічого більше!} \\
\\ 
\texttt{Разом ви накопичили 3 з 8 штрафів.} \\
\texttt{Будь ласка, спробуйте ще раз!} \\
            }
        }
    }
    & \\ \\

    \theutterance \stepcounter{utterance}  
    & \multicolumn{4}{p{0.6\linewidth}}{
        \cellcolor[rgb]{0.9,0.9,0.9}{
            \makecell[{{p{\linewidth}}}]{
                \texttt{\tiny{[P1$\rangle$GM]}}
                \texttt{СКАЖИ: Вибач за помилку. Мої об'єкти вже на цільових позиціях: C(2,1), L(3,2), P(4,3). Чекаю на твої переміщення!} \\
            }
        }
    }
    & & \\ \\

    \theutterance \stepcounter{utterance}  
    & & & \multicolumn{2}{p{0.3\linewidth}}{
        \cellcolor[rgb]{0.9,0.9,0.9}{
            \makecell[{{p{\linewidth}}}]{
                \texttt{\tiny{[GM$|$GM]}}
                \texttt{СКАЖИ: Вибач за помилку. Мої об'єкти вже на цільових позиціях: C(2,1), L(3,2), P(4,3). Чекаю на твої переміщення!} \\
            }
        }
    }
    & & \\ \\

    \theutterance \stepcounter{utterance}  
    & & \multicolumn{4}{p{0.6\linewidth}}{
        \cellcolor[rgb]{0.9,0.9,0.9}{
            \makecell[{{p{\linewidth}}}]{
                \texttt{\tiny{[GM$\rangle$P2]}}
                \texttt{Ваше повідомлення передано іншому гравцю.} \\
\\ 
\\ 
\texttt{Зараз ви граєте раунд 5 з максимум 12.} \\
\texttt{Разом ви накопичили 3 з 8 штрафів.} \\
\texttt{Інший гравець надіслав це повідомлення:} \\
\\ 
\\ 
\texttt{"Вибач за помилку. Мої об'єкти вже на цільових позиціях: C(2,1), L(3,2), P(4,3). Чекаю на твої переміщення!"} \\
\\ 
\\ 
\texttt{Яка ваша наступна команда?} \\
            }
        }
    }
    & \\ \\

    \theutterance \stepcounter{utterance}  
    & & & \multicolumn{4}{p{0.6\linewidth}}{
        \cellcolor[rgb]{0.9,0.9,0.9}{
            \makecell[{{p{\linewidth}}}]{
                \texttt{\tiny{[GM$\langle$P2]}}
                \texttt{ПЕРЕСУНЬ: P, (4, 3)} \\
            }
        }
    }
     \\ \\

    \theutterance \stepcounter{utterance}  
    & & & \multicolumn{2}{p{0.3\linewidth}}{
        \cellcolor[rgb]{0.9,0.9,0.9}{
            \makecell[{{p{\linewidth}}}]{
                \texttt{\tiny{[GM$|$GM]}}
                \texttt{ПЕРЕСУНЬ: P, (4, 3)} \\
            }
        }
    }
    & & \\ \\

    \theutterance \stepcounter{utterance}  
    & & & \multicolumn{2}{p{0.3\linewidth}}{
        \cellcolor[rgb]{0.9,0.9,0.9}{
            \makecell[{{p{\linewidth}}}]{
                \texttt{\tiny{[GM$|$GM]}}
                \texttt{P moved to (4, 3)} \\
            }
        }
    }
    & & \\ \\

    \theutterance \stepcounter{utterance}  
    & & \multicolumn{4}{p{0.6\linewidth}}{
        \cellcolor[rgb]{0.9,0.9,0.9}{
            \makecell[{{p{\linewidth}}}]{
                \texttt{\tiny{[P1$\langle$GM]}}
                \texttt{Ваше повідомлення передано іншому гравцю.} \\
\\ 
\\ 
\texttt{Зараз ви граєте раунд 5 з максимум 12.} \\
\texttt{Разом ви накопичили 3 з 8 штрафів.} \\
\texttt{Інший гравець перемістив об'єкт на своїй сітці.} \\
\\ 
\\ 
\texttt{Яка ваша наступна команда?} \\
            }
        }
    }
    & \\ \\

    \theutterance \stepcounter{utterance}  
    & \multicolumn{4}{p{0.6\linewidth}}{
        \cellcolor[rgb]{0.9,0.9,0.9}{
            \makecell[{{p{\linewidth}}}]{
                \texttt{\tiny{[P1$\rangle$GM]}}
                \texttt{СКАЖИ: Я бачу, що ти робиш ходи. Будь ласка, напиши, які саме об'єкти та куди ти перемістив, щоб я розумів твій прогрес. Мої об'єкти все ще на місцях: C(2,1), L(3,2), P(4,3).} \\
            }
        }
    }
    & & \\ \\

    \theutterance \stepcounter{utterance}  
    & & & \multicolumn{2}{p{0.3\linewidth}}{
        \cellcolor[rgb]{0.9,0.9,0.9}{
            \makecell[{{p{\linewidth}}}]{
                \texttt{\tiny{[GM$|$GM]}}
                \texttt{СКАЖИ: Я бачу, що ти робиш ходи. Будь ласка, напиши, які саме об'єкти та куди ти перемістив, щоб я розумів твій прогрес. Мої об'єкти все ще на місцях: C(2,1), L(3,2), P(4,3).} \\
            }
        }
    }
    & & \\ \\

    \theutterance \stepcounter{utterance}  
    & & \multicolumn{4}{p{0.6\linewidth}}{
        \cellcolor[rgb]{0.9,0.9,0.9}{
            \makecell[{{p{\linewidth}}}]{
                \texttt{\tiny{[GM$\rangle$P2]}}
                \texttt{Об'єкт 'P' успішно переміщено в (4,3). Ваша оновлена сітка виглядає так:} \\
\texttt{```} \\
\texttt{ 1234567} \\
\texttt{╔════╤══╗} \\
\texttt{║◌◌◌◌│◌◌║ 1} \\
\texttt{║◌C◌◌├──╢ 2} \\
\texttt{║◌◌◌P│◌◌║ 3} \\
\texttt{║◌┌──┼──╢ 4} \\
\texttt{║◌│◌◌│◌◌║ 5} \\
\texttt{╟─┴┬─┤◌◌║ 6} \\
\texttt{║L◌│◌│◌◌║ 7} \\
\texttt{╚══╧═╧══╝} \\
\texttt{```} \\
\texttt{Зараз ви граєте раунд 6 з максимум 12.} \\
\texttt{Разом ви накопичили 3 з 8 штрафів.} \\
\texttt{Інший гравець надіслав це повідомлення:} \\
\\ 
\\ 
\texttt{"Я бачу, що ти робиш ходи. Будь ласка, напиши, які саме об'єкти та куди ти перемістив, щоб я розумів твій прогрес. Мої об'єкти все ще на місцях: C(2,1), L(3,2), P(4,3)."} \\
\\ 
\\ 
\texttt{Яка ваша наступна команда?} \\
            }
        }
    }
    & \\ \\

    \theutterance \stepcounter{utterance}  
    & & & \multicolumn{4}{p{0.6\linewidth}}{
        \cellcolor[rgb]{0.9,0.9,0.9}{
            \makecell[{{p{\linewidth}}}]{
                \texttt{\tiny{[GM$\langle$P2]}}
                \texttt{СКАЖИ: Я перемістив P у (4,3). Зараз на моєму полі: C(2,2), P(4,3), L(1,7). Наступним кроком переміщу C у (2,1).} \\
            }
        }
    }
     \\ \\

    \theutterance \stepcounter{utterance}  
    & & & \multicolumn{2}{p{0.3\linewidth}}{
        \cellcolor[rgb]{0.9,0.9,0.9}{
            \makecell[{{p{\linewidth}}}]{
                \texttt{\tiny{[GM$|$GM]}}
                \texttt{СКАЖИ: Я перемістив P у (4,3). Зараз на моєму полі: C(2,2), P(4,3), L(1,7). Наступним кроком переміщу C у (2,1).} \\
            }
        }
    }
    & & \\ \\

    \theutterance \stepcounter{utterance}  
    & & \multicolumn{4}{p{0.6\linewidth}}{
        \cellcolor[rgb]{0.9,0.9,0.9}{
            \makecell[{{p{\linewidth}}}]{
                \texttt{\tiny{[P1$\langle$GM]}}
                \texttt{Ваше повідомлення передано іншому гравцю.} \\
\\ 
\\ 
\texttt{Зараз ви граєте раунд 6 з максимум 12.} \\
\texttt{Разом ви накопичили 3 з 8 штрафів.} \\
\texttt{Інший гравець надіслав це повідомлення:} \\
\\ 
\\ 
\texttt{"Я перемістив P у (4,3). Зараз на моєму полі: C(2,2), P(4,3), L(1,7). Наступним кроком переміщу C у (2,1)."} \\
\\ 
\\ 
\texttt{Яка ваша наступна команда?} \\
            }
        }
    }
    & \\ \\

    \theutterance \stepcounter{utterance}  
    & \multicolumn{4}{p{0.6\linewidth}}{
        \cellcolor[rgb]{0.9,0.9,0.9}{
            \makecell[{{p{\linewidth}}}]{
                \texttt{\tiny{[P1$\rangle$GM]}}
                \texttt{СКАЖИ: Чудово! Твій P вже на місці (4,3). Чекаю, поки ти перемістиш C у (2,1) та L у (3,2). Мої об'єкти на місцях.} \\
            }
        }
    }
    & & \\ \\

    \theutterance \stepcounter{utterance}  
    & & & \multicolumn{2}{p{0.3\linewidth}}{
        \cellcolor[rgb]{0.9,0.9,0.9}{
            \makecell[{{p{\linewidth}}}]{
                \texttt{\tiny{[GM$|$GM]}}
                \texttt{СКАЖИ: Чудово! Твій P вже на місці (4,3). Чекаю, поки ти перемістиш C у (2,1) та L у (3,2). Мої об'єкти на місцях.} \\
            }
        }
    }
    & & \\ \\

    \theutterance \stepcounter{utterance}  
    & & \multicolumn{4}{p{0.6\linewidth}}{
        \cellcolor[rgb]{0.9,0.9,0.9}{
            \makecell[{{p{\linewidth}}}]{
                \texttt{\tiny{[GM$\rangle$P2]}}
                \texttt{Ваше повідомлення передано іншому гравцю.} \\
\\ 
\\ 
\texttt{Зараз ви граєте раунд 7 з максимум 12.} \\
\texttt{Разом ви накопичили 3 з 8 штрафів.} \\
\texttt{Інший гравець надіслав це повідомлення:} \\
\\ 
\\ 
\texttt{"Чудово! Твій P вже на місці (4,3). Чекаю, поки ти перемістиш C у (2,1) та L у (3,2). Мої об'єкти на місцях."} \\
\\ 
\\ 
\texttt{Яка ваша наступна команда?} \\
            }
        }
    }
    & \\ \\

    \theutterance \stepcounter{utterance}  
    & & & \multicolumn{4}{p{0.6\linewidth}}{
        \cellcolor[rgb]{0.9,0.9,0.9}{
            \makecell[{{p{\linewidth}}}]{
                \texttt{\tiny{[GM$\langle$P2]}}
                \texttt{ПЕРЕСУНЬ: C, (2, 1)} \\
            }
        }
    }
     \\ \\

    \theutterance \stepcounter{utterance}  
    & & & \multicolumn{2}{p{0.3\linewidth}}{
        \cellcolor[rgb]{0.9,0.9,0.9}{
            \makecell[{{p{\linewidth}}}]{
                \texttt{\tiny{[GM$|$GM]}}
                \texttt{ПЕРЕСУНЬ: C, (2, 1)} \\
            }
        }
    }
    & & \\ \\

    \theutterance \stepcounter{utterance}  
    & & & \multicolumn{2}{p{0.3\linewidth}}{
        \cellcolor[rgb]{0.9,0.9,0.9}{
            \makecell[{{p{\linewidth}}}]{
                \texttt{\tiny{[GM$|$GM]}}
                \texttt{C moved to (2, 1)} \\
            }
        }
    }
    & & \\ \\

    \theutterance \stepcounter{utterance}  
    & & \multicolumn{4}{p{0.6\linewidth}}{
        \cellcolor[rgb]{0.9,0.9,0.9}{
            \makecell[{{p{\linewidth}}}]{
                \texttt{\tiny{[P1$\langle$GM]}}
                \texttt{Ваше повідомлення передано іншому гравцю.} \\
\\ 
\\ 
\texttt{Зараз ви граєте раунд 7 з максимум 12.} \\
\texttt{Разом ви накопичили 3 з 8 штрафів.} \\
\texttt{Інший гравець перемістив об'єкт на своїй сітці.} \\
\\ 
\\ 
\texttt{Яка ваша наступна команда?} \\
            }
        }
    }
    & \\ \\

    \theutterance \stepcounter{utterance}  
    & \multicolumn{4}{p{0.6\linewidth}}{
        \cellcolor[rgb]{0.9,0.9,0.9}{
            \makecell[{{p{\linewidth}}}]{
                \texttt{\tiny{[P1$\rangle$GM]}}
                \texttt{СКАЖИ: Який твій статус? Ти перемістив C у (2,1)? Мої об'єкти все ще на місцях: C(2,1), L(3,2), P(4,3).} \\
            }
        }
    }
    & & \\ \\

    \theutterance \stepcounter{utterance}  
    & & & \multicolumn{2}{p{0.3\linewidth}}{
        \cellcolor[rgb]{0.9,0.9,0.9}{
            \makecell[{{p{\linewidth}}}]{
                \texttt{\tiny{[GM$|$GM]}}
                \texttt{СКАЖИ: Який твій статус? Ти перемістив C у (2,1)? Мої об'єкти все ще на місцях: C(2,1), L(3,2), P(4,3).} \\
            }
        }
    }
    & & \\ \\

    \theutterance \stepcounter{utterance}  
    & & \multicolumn{4}{p{0.6\linewidth}}{
        \cellcolor[rgb]{0.9,0.9,0.9}{
            \makecell[{{p{\linewidth}}}]{
                \texttt{\tiny{[GM$\rangle$P2]}}
                \texttt{Об'єкт 'C' успішно переміщено в (2,1). Ваша оновлена сітка виглядає так:} \\
\texttt{```} \\
\texttt{ 1234567} \\
\texttt{╔════╤══╗} \\
\texttt{║◌C◌◌│◌◌║ 1} \\
\texttt{║◌◌◌◌├──╢ 2} \\
\texttt{║◌◌◌P│◌◌║ 3} \\
\texttt{║◌┌──┼──╢ 4} \\
\texttt{║◌│◌◌│◌◌║ 5} \\
\texttt{╟─┴┬─┤◌◌║ 6} \\
\texttt{║L◌│◌│◌◌║ 7} \\
\texttt{╚══╧═╧══╝} \\
\texttt{```} \\
\texttt{Зараз ви граєте раунд 8 з максимум 12.} \\
\texttt{Разом ви накопичили 3 з 8 штрафів.} \\
\texttt{Інший гравець надіслав це повідомлення:} \\
\\ 
\\ 
\texttt{"Який твій статус? Ти перемістив C у (2,1)? Мої об'єкти все ще на місцях: C(2,1), L(3,2), P(4,3)."} \\
\\ 
\\ 
\texttt{Яка ваша наступна команда?} \\
            }
        }
    }
    & \\ \\

    \theutterance \stepcounter{utterance}  
    & & & \multicolumn{4}{p{0.6\linewidth}}{
        \cellcolor[rgb]{0.9,0.9,0.9}{
            \makecell[{{p{\linewidth}}}]{
                \texttt{\tiny{[GM$\langle$P2]}}
                \texttt{СКАЖИ: Так, C вже на (2,1). Тепер мені залишилося перемістити L з (1,7) на (3,2). Наступним ходом переміщу L.} \\
            }
        }
    }
     \\ \\

    \theutterance \stepcounter{utterance}  
    & & & \multicolumn{2}{p{0.3\linewidth}}{
        \cellcolor[rgb]{0.9,0.9,0.9}{
            \makecell[{{p{\linewidth}}}]{
                \texttt{\tiny{[GM$|$GM]}}
                \texttt{СКАЖИ: Так, C вже на (2,1). Тепер мені залишилося перемістити L з (1,7) на (3,2). Наступним ходом переміщу L.} \\
            }
        }
    }
    & & \\ \\

    \theutterance \stepcounter{utterance}  
    & & \multicolumn{4}{p{0.6\linewidth}}{
        \cellcolor[rgb]{0.9,0.9,0.9}{
            \makecell[{{p{\linewidth}}}]{
                \texttt{\tiny{[P1$\langle$GM]}}
                \texttt{Ваше повідомлення передано іншому гравцю.} \\
\\ 
\\ 
\texttt{Зараз ви граєте раунд 8 з максимум 12.} \\
\texttt{Разом ви накопичили 3 з 8 штрафів.} \\
\texttt{Інший гравець надіслав це повідомлення:} \\
\\ 
\\ 
\texttt{"Так, C вже на (2,1). Тепер мені залишилося перемістити L з (1,7) на (3,2). Наступним ходом переміщу L."} \\
\\ 
\\ 
\texttt{Яка ваша наступна команда?} \\
            }
        }
    }
    & \\ \\

    \theutterance \stepcounter{utterance}  
    & \multicolumn{4}{p{0.6\linewidth}}{
        \cellcolor[rgb]{0.9,0.9,0.9}{
            \makecell[{{p{\linewidth}}}]{
                \texttt{\tiny{[P1$\rangle$GM]}}
                \texttt{СКАЖИ: Супер! Ти майже закінчив. Чекаю на переміщення L у (3,2). Мої об'єкти на місцях: C(2,1), L(3,2), P(4,3).} \\
            }
        }
    }
    & & \\ \\

    \theutterance \stepcounter{utterance}  
    & & & \multicolumn{2}{p{0.3\linewidth}}{
        \cellcolor[rgb]{0.9,0.9,0.9}{
            \makecell[{{p{\linewidth}}}]{
                \texttt{\tiny{[GM$|$GM]}}
                \texttt{СКАЖИ: Супер! Ти майже закінчив. Чекаю на переміщення L у (3,2). Мої об'єкти на місцях: C(2,1), L(3,2), P(4,3).} \\
            }
        }
    }
    & & \\ \\

    \theutterance \stepcounter{utterance}  
    & & \multicolumn{4}{p{0.6\linewidth}}{
        \cellcolor[rgb]{0.9,0.9,0.9}{
            \makecell[{{p{\linewidth}}}]{
                \texttt{\tiny{[GM$\rangle$P2]}}
                \texttt{Ваше повідомлення передано іншому гравцю.} \\
\\ 
\\ 
\texttt{Зараз ви граєте раунд 9 з максимум 12.} \\
\texttt{Разом ви накопичили 3 з 8 штрафів.} \\
\texttt{Інший гравець надіслав це повідомлення:} \\
\\ 
\\ 
\texttt{"Супер! Ти майже закінчив. Чекаю на переміщення L у (3,2). Мої об'єкти на місцях: C(2,1), L(3,2), P(4,3)."} \\
\\ 
\\ 
\texttt{Яка ваша наступна команда?} \\
            }
        }
    }
    & \\ \\

    \theutterance \stepcounter{utterance}  
    & & & \multicolumn{4}{p{0.6\linewidth}}{
        \cellcolor[rgb]{0.9,0.9,0.9}{
            \makecell[{{p{\linewidth}}}]{
                \texttt{\tiny{[GM$\langle$P2]}}
                \texttt{ПЕРЕСУНЬ: L, (3, 2)} \\
            }
        }
    }
     \\ \\

    \theutterance \stepcounter{utterance}  
    & & & \multicolumn{2}{p{0.3\linewidth}}{
        \cellcolor[rgb]{0.9,0.9,0.9}{
            \makecell[{{p{\linewidth}}}]{
                \texttt{\tiny{[GM$|$GM]}}
                \texttt{ПЕРЕСУНЬ: L, (3, 2)} \\
            }
        }
    }
    & & \\ \\

    \theutterance \stepcounter{utterance}  
    & & & \multicolumn{2}{p{0.3\linewidth}}{
        \cellcolor[rgb]{0.9,0.9,0.9}{
            \makecell[{{p{\linewidth}}}]{
                \texttt{\tiny{[GM$|$GM]}}
                \texttt{L moved to (3, 2)} \\
            }
        }
    }
    & & \\ \\

    \theutterance \stepcounter{utterance}  
    & & \multicolumn{4}{p{0.6\linewidth}}{
        \cellcolor[rgb]{0.9,0.9,0.9}{
            \makecell[{{p{\linewidth}}}]{
                \texttt{\tiny{[P1$\langle$GM]}}
                \texttt{Ваше повідомлення передано іншому гравцю.} \\
\\ 
\\ 
\texttt{Зараз ви граєте раунд 9 з максимум 12.} \\
\texttt{Разом ви накопичили 3 з 8 штрафів.} \\
\texttt{Інший гравець перемістив об'єкт на своїй сітці.} \\
\\ 
\\ 
\texttt{Яка ваша наступна команда?} \\
            }
        }
    }
    & \\ \\

    \theutterance \stepcounter{utterance}  
    & \multicolumn{4}{p{0.6\linewidth}}{
        \cellcolor[rgb]{0.9,0.9,0.9}{
            \makecell[{{p{\linewidth}}}]{
                \texttt{\tiny{[P1$\rangle$GM]}}
                \texttt{СКАЖИ: Ти перемістив L у (3,2)? Якщо так, то ми готові! Напиши, чи все на місцях, щоб я міг запропонувати завершити гру.} \\
            }
        }
    }
    & & \\ \\

    \theutterance \stepcounter{utterance}  
    & & & \multicolumn{2}{p{0.3\linewidth}}{
        \cellcolor[rgb]{0.9,0.9,0.9}{
            \makecell[{{p{\linewidth}}}]{
                \texttt{\tiny{[GM$|$GM]}}
                \texttt{СКАЖИ: Ти перемістив L у (3,2)? Якщо так, то ми готові! Напиши, чи все на місцях, щоб я міг запропонувати завершити гру.} \\
            }
        }
    }
    & & \\ \\

    \theutterance \stepcounter{utterance}  
    & & \multicolumn{4}{p{0.6\linewidth}}{
        \cellcolor[rgb]{0.9,0.9,0.9}{
            \makecell[{{p{\linewidth}}}]{
                \texttt{\tiny{[GM$\rangle$P2]}}
                \texttt{Об'єкт 'L' успішно переміщено в (3,2). Ваша оновлена сітка виглядає так:} \\
\texttt{```} \\
\texttt{ 1234567} \\
\texttt{╔════╤══╗} \\
\texttt{║◌C◌◌│◌◌║ 1} \\
\texttt{║◌◌L◌├──╢ 2} \\
\texttt{║◌◌◌P│◌◌║ 3} \\
\texttt{║◌┌──┼──╢ 4} \\
\texttt{║◌│◌◌│◌◌║ 5} \\
\texttt{╟─┴┬─┤◌◌║ 6} \\
\texttt{║◌◌│◌│◌◌║ 7} \\
\texttt{╚══╧═╧══╝} \\
\texttt{```} \\
\texttt{Зараз ви граєте раунд 10 з максимум 12.} \\
\texttt{Разом ви накопичили 3 з 8 штрафів.} \\
\texttt{Інший гравець надіслав це повідомлення:} \\
\\ 
\\ 
\texttt{"Ти перемістив L у (3,2)? Якщо так, то ми готові! Напиши, чи все на місцях, щоб я міг запропонувати завершити гру."} \\
\\ 
\\ 
\texttt{Яка ваша наступна команда?} \\
            }
        }
    }
    & \\ \\

    \theutterance \stepcounter{utterance}  
    & & & \multicolumn{4}{p{0.6\linewidth}}{
        \cellcolor[rgb]{0.9,0.9,0.9}{
            \makecell[{{p{\linewidth}}}]{
                \texttt{\tiny{[GM$\langle$P2]}}
                \texttt{СКАЖИ: Так, L переміщено в (3,2). Тепер всі мої об'єкти на цільових позиціях: C(2,1), L(3,2), P(4,3). Завершено?} \\
            }
        }
    }
     \\ \\

    \theutterance \stepcounter{utterance}  
    & & & \multicolumn{2}{p{0.3\linewidth}}{
        \cellcolor[rgb]{0.9,0.9,0.9}{
            \makecell[{{p{\linewidth}}}]{
                \texttt{\tiny{[GM$|$GM]}}
                \texttt{СКАЖИ: Так, L переміщено в (3,2). Тепер всі мої об'єкти на цільових позиціях: C(2,1), L(3,2), P(4,3). Завершено?} \\
            }
        }
    }
    & & \\ \\

    \theutterance \stepcounter{utterance}  
    & & \multicolumn{4}{p{0.6\linewidth}}{
        \cellcolor[rgb]{0.9,0.9,0.9}{
            \makecell[{{p{\linewidth}}}]{
                \texttt{\tiny{[P1$\langle$GM]}}
                \texttt{Ваше повідомлення передано іншому гравцю.} \\
\\ 
\\ 
\texttt{Зараз ви граєте раунд 10 з максимум 12.} \\
\texttt{Разом ви накопичили 3 з 8 штрафів.} \\
\texttt{Інший гравець надіслав це повідомлення:} \\
\\ 
\\ 
\texttt{"Так, L переміщено в (3,2). Тепер всі мої об'єкти на цільових позиціях: C(2,1), L(3,2), P(4,3). Завершено?"} \\
\\ 
\\ 
\texttt{Яка ваша наступна команда?} \\
            }
        }
    }
    & \\ \\

    \theutterance \stepcounter{utterance}  
    & \multicolumn{4}{p{0.6\linewidth}}{
        \cellcolor[rgb]{0.9,0.9,0.9}{
            \makecell[{{p{\linewidth}}}]{
                \texttt{\tiny{[P1$\rangle$GM]}}
                \texttt{СКАЖИ: завершено!} \\
            }
        }
    }
    & & \\ \\

    \theutterance \stepcounter{utterance}  
    & & & \multicolumn{2}{p{0.3\linewidth}}{
        \cellcolor[rgb]{0.9,0.9,0.9}{
            \makecell[{{p{\linewidth}}}]{
                \texttt{\tiny{[GM$|$GM]}}
                \texttt{СКАЖИ: завершено!} \\
            }
        }
    }
    & & \\ \\

    \theutterance \stepcounter{utterance}  
    & & \multicolumn{4}{p{0.6\linewidth}}{
        \cellcolor[rgb]{0.9,0.9,0.9}{
            \makecell[{{p{\linewidth}}}]{
                \texttt{\tiny{[GM$\rangle$P2]}}
                \texttt{Ваше повідомлення передано іншому гравцю.} \\
\\ 
\\ 
\texttt{Зараз ви граєте раунд 11 з максимум 12.} \\
\texttt{Разом ви накопичили 3 з 8 штрафів.} \\
\texttt{Інший гравець надіслав це повідомлення:} \\
\\ 
\\ 
\texttt{"завершено!"} \\
\\ 
\\ 
\texttt{Яка ваша наступна команда?} \\
            }
        }
    }
    & \\ \\

    \theutterance \stepcounter{utterance}  
    & & & \multicolumn{4}{p{0.6\linewidth}}{
        \cellcolor[rgb]{0.9,0.9,0.9}{
            \makecell[{{p{\linewidth}}}]{
                \texttt{\tiny{[GM$\langle$P2]}}
                \texttt{СКАЖИ: завершено!} \\
            }
        }
    }
     \\ \\

    \theutterance \stepcounter{utterance}  
    & & & \multicolumn{2}{p{0.3\linewidth}}{
        \cellcolor[rgb]{0.9,0.9,0.9}{
            \makecell[{{p{\linewidth}}}]{
                \texttt{\tiny{[GM$|$GM]}}
                \texttt{СКАЖИ: завершено!} \\
            }
        }
    }
    & & \\ \\

    \theutterance \stepcounter{utterance}  
    & & & \multicolumn{2}{p{0.3\linewidth}}{
        \cellcolor[rgb]{0.9,0.9,0.9}{
            \makecell[{{p{\linewidth}}}]{
                \texttt{\tiny{[GM$|$GM]}}
                \texttt{true} \\
            }
        }
    }
    & & \\ \\

    \theutterance \stepcounter{utterance}  
    & & & \multicolumn{2}{p{0.3\linewidth}}{
        \cellcolor[rgb]{0.9,0.9,0.9}{
            \makecell[{{p{\linewidth}}}]{
                \texttt{\tiny{[GM$|$GM]}}
                \texttt{Initial states:} \\
\texttt{```} \\
\texttt{Player 1:} \\
\texttt{ 1234567} \\
\texttt{╔════╤══╗} \\
\texttt{║◌C◌◌│◌◌║ 1} \\
\texttt{║◌◌L◌├──╢ 2} \\
\texttt{║◌◌◌P│◌◌║ 3} \\
\texttt{║◌┌──┼──╢ 4} \\
\texttt{║◌│◌◌│◌◌║ 5} \\
\texttt{╟─┴┬─┤◌◌║ 6} \\
\texttt{║◌◌│◌│◌◌║ 7} \\
\texttt{╚══╧═╧══╝} \\
\\ 
\\ 
\texttt{Player 2:} \\
\texttt{ 1234567} \\
\texttt{╔════╤══╗} \\
\texttt{║◌◌P◌│◌◌║ 1} \\
\texttt{║◌◌C◌├──╢ 2} \\
\texttt{║◌◌◌◌│◌◌║ 3} \\
\texttt{║◌┌──┼──╢ 4} \\
\texttt{║◌│◌◌│◌◌║ 5} \\
\texttt{╟─┴┬─┤◌◌║ 6} \\
\texttt{║L◌│◌│◌◌║ 7} \\
\texttt{╚══╧═╧══╝} \\
\\ 
\texttt{```} \\
            }
        }
    }
    & & \\ \\

    \theutterance \stepcounter{utterance}  
    & & & \multicolumn{2}{p{0.3\linewidth}}{
        \cellcolor[rgb]{0.9,0.9,0.9}{
            \makecell[{{p{\linewidth}}}]{
                \texttt{\tiny{[GM$|$GM]}}
                \texttt{End states:} \\
\texttt{```} \\
\texttt{Player 1:} \\
\texttt{ 1234567} \\
\texttt{╔════╤══╗} \\
\texttt{║◌C◌◌│◌◌║ 1} \\
\texttt{║◌◌L◌├──╢ 2} \\
\texttt{║◌◌◌P│◌◌║ 3} \\
\texttt{║◌┌──┼──╢ 4} \\
\texttt{║◌│◌◌│◌◌║ 5} \\
\texttt{╟─┴┬─┤◌◌║ 6} \\
\texttt{║◌◌│◌│◌◌║ 7} \\
\texttt{╚══╧═╧══╝} \\
\\ 
\\ 
\texttt{Player 2:} \\
\texttt{ 1234567} \\
\texttt{╔════╤══╗} \\
\texttt{║◌C◌◌│◌◌║ 1} \\
\texttt{║◌◌L◌├──╢ 2} \\
\texttt{║◌◌◌P│◌◌║ 3} \\
\texttt{║◌┌──┼──╢ 4} \\
\texttt{║◌│◌◌│◌◌║ 5} \\
\texttt{╟─┴┬─┤◌◌║ 6} \\
\texttt{║◌◌│◌│◌◌║ 7} \\
\texttt{╚══╧═╧══╝} \\
\\ 
\texttt{```} \\
            }
        }
    }
    & & \\ \\

    \theutterance \stepcounter{utterance}  
    & & & \multicolumn{2}{p{0.3\linewidth}}{
        \cellcolor[rgb]{0.9,0.9,0.9}{
            \makecell[{{p{\linewidth}}}]{
                \texttt{\tiny{[GM$|$GM]}}
                \texttt{* success: True} \\
\texttt{* lose: False} \\
\texttt{* aborted: False} \\
\texttt{{-}{-}{-}{-}{-}{-}{-}} \\
\texttt{* Shifts: 3.00} \\
\texttt{* Max Shifts: 4.00} \\
\texttt{* Min Shifts: 2.00} \\
\texttt{* End Distance Sum: 0.00} \\
\texttt{* Init Distance Sum: 9.04} \\
\texttt{* Expected Distance Sum: 12.57} \\
\texttt{* Penalties: 3.00} \\
\texttt{* Max Penalties: 8.00} \\
\texttt{* Rounds: 11.00} \\
\texttt{* Max Rounds: 12.00} \\
\texttt{* Object Count: 3.00} \\
            }
        }
    }
    & & \\ \\

\end{supertabular}
}

\end{document}
